\documentclass[../BTL.tex]{subfiles}
\begin{document}

\section{Đặt vấn đề}
\label{section:1.1}

Bài toán tìm đường đi ngắn nhất trên lưới hai chiều phát sinh từ nhiều ứng dụng thực tiễn. Trong robot di động tự hành, khả năng xác định hành trình tối ưu từ vị trí xuất phát đến điểm đích trên bản đồ có chướng ngại vật là yêu cầu nền tảng cho điều khiển và tránh va chạm. Các hệ thống trò chơi điện tử chiến thuật và nhập vai trực tuyến cũng yêu cầu hàng nghìn thực thể ảo đồng thời tính toán đường đi với chi phí xử lý thấp. Lĩnh vực logistics kho bãi và hệ thống giao thông thông minh đối mặt với bài toán tương tự trên không gian lưới được số hóa.

Thuật toán A* do Hart, Nilsson và Raphael công bố trên IEEE Transactions năm 1968 đặt nền tảng lý thuyết cho bài toán này. Nguyên lý cốt lõi của A* là hàm đánh giá \(f(n) = g(n) + h(n)\), trong đó \(g(n)\) là chi phí thực tế từ điểm xuất phát đến đỉnh \(n\), \(h(n)\) là chi phí ước lượng từ \(n\) đến đích. Chứng minh tính tối ưu của A* khi \(h(n)\) là hàm chấp nhận được tạo ra bước đột phá so với các thuật toán tìm kiếm mù như Dijkstra hay BFS. Hiệu năng của A* trong không gian trạng thái lớn phụ thuộc vào ba yếu tố: lựa chọn hàm heuristic, khả năng di chuyển theo đường chéo, và chiến lược tìm kiếm đơn hướng so với song hướng.

Khi bài toán tìm đường được giải quyết hiệu quả, các hệ thống ứng dụng thu được ba lợi ích chính. Robot giảm thời gian di chuyển và tiết kiệm năng lượng tiêu thụ trong quá trình hoạt động. Trò chơi điện tử tăng số lượng thực thể có thể xử lý đồng thời trên cùng nền tảng phần cứng. Hệ thống thời gian thực đáp ứng được các ràng buộc về độ trễ tính toán. Các giải pháp này còn có thể mở rộng sang lĩnh vực lập quỹ đạo cho cánh tay robot trong không gian làm việc nhiều chiều, hoặc tối ưu hóa lộ trình vận chuyển trong kho tự động hóa. Phần tiếp theo sẽ khảo sát các nghiên cứu hiện có để xác định mục tiêu cụ thể cho đề tài.

\section{Mục tiêu và phạm vi đề tài}
\label{section:1.2}

Các nghiên cứu cải tiến A* trên lưới hai chiều tập trung vào ba hướng chính. Hướng thứ nhất là điều chỉnh hàm heuristic với bốn biến thể phổ biến: Manhattan (\(|dx|+|dy|\)), Euclidean (\(\sqrt{dx^2+dy^2}\)), Octile (\(\max(dx,dy)+(\sqrt{2}-1)\times\min(dx,dy)\)), Chebyshev (\(\max(dx,dy)\)). Trường hợp đặc biệt Dijkstra (\(h=0\)) biến A* thành thuật toán tìm kiếm không có thông tin heuristic. Hướng thứ hai là cho phép di chuyển góc bất kỳ, tiêu biểu là Theta* do Nash và cộng sự giới thiệu tại AAAI 2007, loại bỏ ràng buộc 8 hướng của lưới truyền thống. Hướng thứ ba là tăng tốc tìm kiếm với Jump Point Search (JPS) của Harabor và Grastien tại AAAI 2011, bỏ qua các đỉnh trung gian không cần thiết trên lưới đồng nhất.

\begin{table}[h]
\centering
\caption{So sánh các phương pháp tìm đường tiêu biểu}
\label{table:1.1}
\begin{tabular}{|l|c|c|c|}
\hline
\textbf{Phương pháp} & \textbf{Độ phức tạp thời gian} & \textbf{Chất lượng đường đi} & \textbf{Khả năng tùy chỉnh} \\
\hline
A* Manhattan & \(O(b^d)\) trung bình & Tối ưu 8 hướng & Cao (4 heuristic) \\
\hline
A* Euclidean & \(O(b^d)\) trung bình & Tối ưu góc bất kỳ & Cao (4 heuristic) \\
\hline
Theta* & \(O(b^d)\) cao hơn & Tối ưu góc bất kỳ & Trung bình \\
\hline
JPS & \(O(n)\) trên lưới đồng nhất & Tối ưu 8 hướng & Thấp (không heuristic) \\
\hline
\end{tabular}
\end{table}

Phân tích bảng 1.1 cho thấy mỗi phương pháp chỉ tối ưu cho một tập con điều kiện đầu vào cụ thể. A* Manhattan đơn giản và hiệu quả nhưng chỉ tạo đường đi theo 8 hướng. A* Euclidean tạo đường đi tự nhiên hơn nhưng tốn thời gian tính toán căn bậc hai. Theta* tối ưu về độ dài đường đi nhưng chi phí kiểm tra điểm cao. JPS nhanh nhất trên lưới đồng nhất nhưng không hỗ trợ heuristic tùy ý.

Hạn chế lớn nhất của các hệ thống hiện tại là thiếu một khung làm việc tích hợp cho phép so sánh trực quan các biến thể. Người dùng không thể dễ dàng thay đổi heuristic, chế độ đa hướng, hay chiến lược tìm kiếm song hướng trên cùng một nền tảng. Hầu hết các công bố chỉ cung cấp kết quả mô phỏng trên dữ liệu tổng hợp mà không đi kèm công cụ trực quan hóa quá trình mở rộng đỉnh theo thời gian thực.

Trên cơ sở các hạn chế đã chỉ ra, mục tiêu của bài tập lớn này là phát triển một hệ thống phần mềm trực quan hóa và so sánh các biến thể A* trên lưới hai chiều. Hệ thống cần đáp ứng bốn yêu cầu chức năng chính. Cài đặt đầy đủ A* với năm hàm heuristic (Euclidean, Manhattan, Octile, Chebyshev, Dijkstra). Hỗ trợ tùy chọn di chuyển đường chéo và kiểm soát cắt góc theo nguyên tắc không đi xuyên qua góc chướng ngại vật. Cung cấp cơ chế điều chỉnh tốc độ hiển thị. Cho phép lưu và tải cấu hình bản đồ dưới định dạng JSON để tái lập thực nghiệm.

Điểm khác biệt của hệ thống so với các công bố hiện có nằm ở những tính năng mang tính ứng dụng cao. Đầu tiên là hệ thống trực quan hóa toàn bộ quá trình tìm kiếm với mã màu phân biệt ba trạng thái của đỉnh: open (màu xanh lá), closed (màu xanh ngọc), path (màu vàng). Thứ hai là hệ thống thống kê chi tiết các chỉ số như thời gian xử lý (ms), số nút đã duyệt (visited nodes), và số phép tính (operations). Cuối cùng là tính năng Save Map và Load Map, cho phép trích xuất và tái tạo chính xác các không gian lưới phức tạp (Bẫy Heuristic, Mê cung) dưới định dạng JSON. Các nội dung chi tiết về cài đặt sẽ được trình bày trong Chương 3.

\section{Định hướng giải pháp}
\label{section:1.3}

Bài tập lớn này đề xuất định hướng giải pháp dựa trên ba quyết định kiến trúc chính. Thứ nhất, sử dụng cơ chế lưu trữ sự kiện (event logging) kết hợp với phát lại có điều khiển, trong đó mỗi bước mở rộng đỉnh được ghi nhận thành cặp (cell\_type, position) vào mảng frames. Thứ hai, tổ chức các tham số heuristic và cấu hình di chuyển (allow\_diagonal, dont\_cross\_corners, diagonal\_cost\_one) thành các biến điều khiển độc lập, cho phép so sánh công bằng giữa các cấu hình. Thứ ba, xây dựng giao diện trực quan với thư viện Tkinter, trong đó mỗi ô lưới được biểu diễn bằng một đối tượng hình chữ nhật độc lập để thay đổi màu sắc thời gian thực.

Giải pháp cụ thể được tổ chức trong một lớp App duy nhất quản lý toàn bộ dữ liệu và điều khiển. Ma trận grid\_data kích thước ROWS×COLUMNS lưu trạng thái các ô (0: đường đi, 1: tường). Cấu trúc dữ liệu tìm kiếm sử dụng dictionary g\_scores để lưu chi phí đường đi đến từng ô, set closed\_set để đánh dấu ô đã mở rộng, và hàng đợi ưu tiên open\_set (heapq) để quản lý các ô đang chờ xét. Hàm \_calculate\_heuristic đóng gói năm phương pháp heuristic: Euclidean, Manhattan, Octile, Chebyshev, và Dijkstra (\(h=0\)).

Hàm tìm kiếm chính \_run\_astar\_algorithm được triển khai theo lưu đồ tuần tự truyền thống nhưng có ghi nhận sự kiện. Mỗi lần một ô được thêm vào open set, hàm ghi ('open', pos) vào mảng frames. Mỗi lần một ô được lấy ra khỏi open set để mở rộng, hàm ghi ('closed', pos) vào mảng frames. Hàm \_get\_neighbors xử lý ba cấp độ kiểm soát di chuyển: Allow Diagonal (bật/tắt di chuyển chéo), Don't Cross Corners (kiểm tra hai ô kề cạnh trước khi cho phép đi chéo), Diagonal Cost = 1 (gán chi phí chéo bằng 1.0 thay vì \(\sqrt{2}\)). Quá trình phát lại sử dụng phương thức after() của Tkinter để tạo hiệu ứng hoạt hình tuần tự với độ trễ tùy chỉnh.

Đóng góp chính của bài tập lớn là ba kết quả có thể đo lường và tái tạo. Thứ nhất, một hệ thống phần mềm hoàn chỉnh cho phép người dùng tương tác xây dựng bản đồ bằng chuột (thêm/xóa tường, kéo thả điểm đầu/cuối), lưu trữ và chia sẻ bản đồ bằng JSON (Save/Load Map), thay đổi cấu hình đa dạng (năm heuristic, ba chế độ di chuyển, tốc độ), và quan sát quá trình tìm kiếm với mã màu trực quan. Thứ hai, chức năng Compare All thực hiện lần lượt năm heuristic trên cùng một bản đồ, thu thập bốn chỉ số (Path Cost, Visited Nodes, Execute Time, Operations) và hiển thị kết quả dạng bảng so sánh. Thứ ba, bộ kiểm thử tự động với cấu hình lưu trong file test\_cases.json, đảm bảo tính đúng đắn của thuật toán dưới mọi kịch bản môi trường phức tạp.

\section{Bố cục bài tập lớn}
\label{section:1.4}

Phần còn lại của báo cáo được tổ chức thành ba chương tiếp theo. Chương 2 trình bày cơ sở lý thuyết của bài toán tìm đường trên đồ thị, bao gồm định nghĩa không gian trạng thái, hàm chi phí, tính chất chấp nhận được của heuristic, và phân tích độ phức tạp thời gian, không gian của A*. Chương 3 giới thiệu toàn bộ quá trình thực nghiệm từ thiết kế kiến trúc phần mềm (cấu trúc lớp App, các biến điều khiển, cơ chế ghi nhận sự kiện) đến triển khai các hàm heuristic, hàm lấy lân cận với ba cấp độ kiểm soát di chuyển, hàm tìm kiếm chính, giao diện người dùng với cơ chế kéo thả, điều khiển tốc độ, quản lý lưu/tải bản đồ, cùng chức năng Compare All và bộ kiểm thử tự động. Chương 4 tổng kết các kết quả đạt được, phân tích những hạn chế còn tồn tại, và đề xuất ba hướng phát triển trong tương lai bao gồm tích hợp tìm kiếm song hướng (bidirectional A*), hỗ trợ xuất bản đồ ra file hình ảnh, và tối ưu hiệu năng bằng cách lưu trữ danh sách neighbor tĩnh.

\subsection*{Tổng quan}

Chương 1 đã đặt bài toán tìm đường ngắn nhất trên lưới hai chiều và chỉ ra sự thiếu vắng một hệ thống trực quan hóa tích hợp nhiều biến thể A*. Nội dung chính của chương bao gồm: phần 1.1 đặt vấn đề và phân tích bối cảnh thực tiễn phát sinh bài toán; phần 1.2 tổng quan các nghiên cứu liên quan và xác định mục tiêu, phạm vi đề tài; phần 1.3 đề xuất định hướng giải pháp công nghệ; phần 1.4 mô tả cấu trúc toàn bộ báo cáo. Chương tiếp theo sẽ trình bày chi tiết cơ sở lý thuyết về hàm heuristic, tính chất tối ưu của A*, và các công trình nghiên cứu liên quan.

\subsection*{Kết chương}

Chương này đã xác định chính xác bài toán tìm đường ngắn nhất trên lưới hai chiều và vị trí của thuật toán A* trong bối cảnh các nghiên cứu hiện đại. Việc phân tích bốn phương pháp (A* Manhattan, A* Euclidean, Theta*, JPS) chỉ ra rằng chưa có giải pháp nào tích hợp đầy đủ bốn tính năng: linh hoạt về heuristic, kiểm soát di chuyển linh hoạt, trực quan hóa thời gian thực, và giao diện tương tác đồ họa. Định hướng giải pháp dựa trên cơ chế lưu trữ sự kiện và phát lại có điều khiển, sử dụng hàng đợi ưu tiên heapq và thư viện Tkinter, được đề xuất để lấp đầy khoảng trống này. Ba đóng góp chính của bài tập lớn là hệ thống trực quan hóa hoàn chỉnh tích hợp tính năng Save/Load Map, chức năng Compare All so sánh năm heuristic, và bộ kiểm thử tự động với dữ liệu linh hoạt. Chương tiếp theo sẽ trình bày chi tiết cơ sở lý thuyết về hàm heuristic, tính chất tối ưu của A*, và các công trình nghiên cứu liên quan.

\end{document}